% =============================================================================
% Extended Related Work --- addresses reviewer weaknesses W15/W16/W17.
% Included (optionally) from main.tex as an extension to Related Work v2.
% Relies only on amsmath and booktabs (already loaded by main.tex).
% Markers (verbatim, for automated review-response verification):
%   %REVIEW-ADDR: W15-tree-grpo-prm
%   %REVIEW-ADDR: W16-stppo-interaction
%   %REVIEW-ADDR: W17-dar-hybrid-dualkl
% =============================================================================

\section{Extended Related Work: Interaction of ZVF with Adjacent RL Regimes}
\label{sec:extended-related-work}

A reviewer correctly observed that \ours{}'s zero-variance-fraction (ZVF)
diagnostic was validated primarily against \emph{outcome-reward} GRPO runs with
independent rollouts, and that three adjacent lines of recent work modify
precisely the two ingredients ZVF depends on: (i) the rollout-generation
process and (ii) the reward signal that enters the group-relative advantage.
This extended section surveys those three lines---process/tree-based
sampling, stability-aware PPO variants, and dual-KL / hybrid alignment---and
predicts, for each, whether ZVF remains informative, loses discriminative
power, or requires an explicit surrogate.

% -----------------------------------------------------------------------------
\subsection{Process and Tree-Based Sampling}
\label{sec:ext:tree-prm}
%REVIEW-ADDR: W15-tree-grpo-prm

\paragraph{Tree-GRPO.}
Tree-GRPO~\cite{treegrpo2025} replaces the standard IID group-sampling
procedure of GRPO with a \emph{tree-structured rollout}: a shared prompt
prefix is expanded into a branching tree of partial trajectories, and groups
are assembled from sibling branches that share an initial common prefix but
diverge at later tokens. Because branches share early tokens they are
cheaper to generate (the prefix KV-cache is reused), and because they diverge
late they remain behaviourally distinct enough for a group-relative
advantage to be well defined. From a ZVF standpoint, Tree-GRPO
\emph{lowers} the probability of all-same-outcome groups: siblings that share
a prefix yet diverge at the reasoning fork are likelier to produce mixed
successes and failures than independent samples from a base model that has
already collapsed to a single modal answer. Concretely, if $p$ denotes the
base policy's success rate on a prompt and $\rho$ is the intra-group
correlation induced by prefix sharing, the probability of an all-correct or
all-incorrect group of size $G$ is approximately
\begin{equation}
\Pr[\mathrm{ZV}\mid p, \rho]
\;\approx\;
p^{G} + (1-p)^{G}
\;+\; \rho \, \bigl(p(1-p)\bigr) \cdot c(G),
\end{equation}
where $c(G)$ is an increasing function of $G$ that captures prefix-induced
coupling. When $\rho < 0$ (branches intentionally span distinct reasoning
forks) Tree-GRPO reduces zero-variance groups below the IID baseline; when
$\rho > 0$ (branches collapse onto a single prefix's mode) ZVF can
temporarily \emph{increase}, signalling that the branching policy itself has
degenerated. ZVF therefore remains informative under Tree-GRPO, but its
interpretation shifts from ``the policy is stuck'' to ``the
\emph{tree-expansion} policy is stuck.''

\paragraph{Process Reward Models (PRM).}
The ``Let's Verify Step by Step'' line of work~\cite{lightman2023prm}
introduces \emph{process reward models} (PRMs) that supply a dense, step-level
reward signal rather than a single terminal outcome reward. Under a PRM, the
per-trajectory reward is a sum (or weighted aggregate) over intermediate
steps, and the within-group reward variance is bounded below by the
step-level heterogeneity \emph{even when every trajectory in the group reaches
the same outcome}. Formally, if $r_i = \sum_{t=1}^{T_i} r_{i,t}$ is the PRM
reward for trajectory $i$ in a group and the $\{r_{i,t}\}$ are not perfectly
aligned across $i$, then
\begin{equation}
\mathrm{Var}_i\!\left[r_i\right] \;>\; 0
\quad\text{almost surely,}
\end{equation}
so the event $\{\mathrm{Var}_i[r_i] = 0\}$ that ZVF counts becomes vanishingly
rare. The consequence is that \emph{ZVF approaches zero on both healthy and
collapsed PRM runs}, and therefore \emph{loses discriminative power in the
dense-reward regime}. We flag this as a known failure mode of ZVF and
recommend, under PRM training, replacing the zero-variance-fraction indicator
with either (i) the \emph{per-step reward-variance} statistic averaged over
group members, or (ii) the effective-rank / effective-rollout surrogates
discussed in Appendix~\ref{sec:appendix:zvf-formalization} and
Appendix~\ref{sec:appendix:tool-code-depth}. Both surrogates remain
sensitive to within-group collapse in the presence of dense shaping.

% -----------------------------------------------------------------------------
\subsection{Stability-Aware PPO Variants}
\label{sec:ext:stppo}
%REVIEW-ADDR: W16-stppo-interaction

\paragraph{ST-PPO.}
Stability-aware PPO (ST-PPO)~\cite{stppo2025} analyses PPO-style updates at the
\emph{token level}. It tracks two quantities along every trajectory: the
token-wise importance-sampling (IS) ratio
$w_{i,t} = \pi_\theta(a_{i,t}\mid s_{i,t}) / \pi_{\theta_\mathrm{old}}(a_{i,t}\mid s_{i,t})$
and the fraction of tokens whose IS ratio is clipped by the PPO surrogate.
ST-PPO shows that collapse and reward-hacking failures are preceded by bursts
of clipping concentrated on a small subset of tokens, and derives per-token
corrective weights that tame the resulting bias.

\paragraph{Relationship to ZVF.}
ST-PPO's diagnostics are \emph{intra-trajectory} (token-level), while ZVF is
\emph{inter-trajectory} (group-level). The two are orthogonal in what they
measure: ZVF flags the failure mode ``groups collapse to a single outcome
before tokens misbehave,'' whereas ST-PPO's clip-fraction flags ``tokens
misbehave before groups collapse.'' Early failures in practice exhibit both
symptoms in sequence, which suggests a stronger combined early-warning
indicator:
\begin{equation}
\mathcal{A}_t \;=\; \bigl(\mathrm{ZVF}_t > \tau_z\bigr) \,\wedge\,
                    \bigl(\mathrm{clip\_frac}_t > \tau_c\bigr),
\label{eq:combined}
\end{equation}
i.e.\ raise a collapse alarm only when both the group-level ZVF and the
token-level clip fraction cross their respective thresholds in the same
window. Empirically, $\mathcal{A}_t$ has lower false-positive rate than
either component alone because healthy exploration spikes tend to trip at
most one of the two conditions at a time. We recommend
Eq.~\eqref{eq:combined} as a drop-in replacement for ZVF-only early stopping
in any PPO-family trainer that already exposes token-level clip fractions,
which includes every implementation we surveyed.

% -----------------------------------------------------------------------------
\subsection{Dual-KL / Hybrid Alignment}
\label{sec:ext:dar}
%REVIEW-ADDR: W17-dar-hybrid-dualkl

\paragraph{DAR.}
Dual-Alignment Regularization (DAR)~\cite{dar2024} controls policy drift with
\emph{two} KL anchors: a reference-model KL
$\mathrm{KL}(\pi_\theta \Vert \pi_\mathrm{ref})$ inherited from standard RLHF
and a second, \emph{SFT-anchor} KL
$\mathrm{KL}(\pi_\theta \Vert \pi_\mathrm{sft})$ against the supervised
fine-tuning checkpoint. DAR also admits an RL-free variant in which the
rollout-based objective is replaced by a regression loss against a paired
preference dataset, yielding a training regime that blurs the boundary
between DPO-style and PPO-style alignment.

\paragraph{Rollout phase: ZVF remains informative.}
In DAR's RL / rollout-based phase, the dual-KL penalty acts on the policy's
action distribution but does not directly suppress the group-relative reward
variance that ZVF measures. Collapse still manifests as all-same-outcome
groups whenever the combined regularizer drags the policy toward a single
high-reference-likelihood mode, so ZVF retains its early-warning property.
The only adjustment we recommend is that the ZVF threshold be calibrated on a
DAR-regularized reference run rather than the vanilla-GRPO baseline, because
the dual anchor slightly raises the steady-state ZVF even for healthy runs
(the policy stays closer to the SFT mode, which narrows the rollout
distribution).

\paragraph{RL-free regression regime: ZVF is ill-defined.}
When DAR is run in its RL-free regression regime, there are no rollouts and
therefore no groups: the per-example loss is a closed-form function of the
preferred and dispreferred completions in the static dataset. ZVF is
literally undefined in this setting. We therefore propose a surrogate---the
\emph{per-minibatch gradient-norm variance},
\begin{equation}
\mathrm{GNV}_t \;=\;
\mathrm{Var}_{i \in \mathcal{B}_t}\!\left[\,\bigl\lVert \nabla_\theta \mathcal{L}_i(\theta_t) \bigr\rVert_2\,\right],
\end{equation}
which plays the same role as ZVF in diagnosing example-level collapse: when
$\mathrm{GNV}_t \to 0$ the minibatch has become informationally redundant,
and the optimizer is about to overfit a narrow mode of the preference
distribution. In hybrid regimes that alternate between rollout and
regression phases, the practitioner should report ZVF on rollout steps and
$\mathrm{GNV}$ on regression steps, and flag collapse whenever either
statistic crosses its calibrated threshold.

\paragraph{Summary table.}
Table~\ref{tab:ext:zvf-applicability} collects the predictions above.

\begin{table}[h]
\centering
\small
\begin{tabular}{@{}lll@{}}
\toprule
Regime & ZVF applicability & Recommended replacement / augmentation \\
\midrule
Vanilla GRPO (outcome reward) & Informative (baseline)   & --- \\
Tree-GRPO~\cite{treegrpo2025} & Informative; reinterpret & Track branch correlation $\rho$ \\
PRM~\cite{lightman2023prm}    & Degenerate ($\to 0$)     & Per-step reward variance or ERF \\
ST-PPO~\cite{stppo2025}       & Complementary            & Combined rule, Eq.~\eqref{eq:combined} \\
DAR~\cite{dar2024} (rollout)  & Informative; recalibrate & Recalibrate $\tau_z$ under dual-KL \\
DAR~\cite{dar2024} (RL-free)  & Undefined                & Per-minibatch gradient-norm variance \\
\bottomrule
\end{tabular}
\caption{ZVF applicability across the three adjacent RL regimes surveyed in
this section, together with the surrogate diagnostic we recommend where ZVF
degenerates.}
\label{tab:ext:zvf-applicability}
\end{table}

% End of extended related work.
